problem:  
Let \( M^n \) be a compact, simply connected Riemannian manifold with **positive sectional curvature** admitting an effective, isometric **cohomogeneity‑one** action by a compact, connected Lie group \( G \).  
Then \( M / G = [0,1] \), with singular orbits \( B_- = G / K_- \) and \( B_+ = G / K_+ \), and principal isotropy \( H = K_- \cap K_+ \).  

Denote by \(\mathrm{codim}(B_\pm) = c_\pm\) the normal dimensions. Assume only the standard topological structure of cohomogeneity‑one:
\[
M \cong D^{c_-}(E_-)\cup D^{c_+}(E_+),
\]
where \(E_\pm \to B_\pm\) are the normal sphere bundles with structure group \(K_\pm/H \subset \mathrm{O}(c_\pm)\).

**Motivating Principle:**  
In positive curvature, Wilking’s connectedness lemma implies that the inclusion of any totally geodesic submanifold of codimension \(\ell\) is \((n-2\ell+1)\)-connected, leading to periodicity of cohomology when \(\ell\) is small.  
Although orbit geometry in cohomogeneity‑one manifolds is more flexible, these inclusion‑map connectivity estimates still constrain the possible topology of \(M\) through the two singular strata \(B_-\) and \(B_+\).  
Known positively curved cohomogeneity‑one examples (such as the Aloff–Wallach spaces and Eschenburg spaces) have relatively large total codimension \(c_- + c_+ \approx n\), whereas symmetric examples (CROSSes with standard actions) often satisfy \(c_- + c_+ \ll n\).

**Open Problem / Conjecture:**  
Does there exist a function \( f(n) \), determined by the connectedness estimates of Wilking‑type, such that if
\[
c_- + c_+ \le f(n),
\]
then \( M \) must have **periodic rational cohomology**, and, in the extremal regime (e.g. \(f(n)\) asymptotic to \(n/2\)), \(M\) must in fact be **G‑equivariantly diffeomorphic to a compact rank‑one symmetric space** with its standard cohomogeneity‑one action?  

In particular:
- Can such a threshold be derived from Wilking’s connectedness lemma applied iteratively to the pair \((B_-, B_+)\)?  
- If so, does it produce a “cohomogeneity‑codimension rigidity line” separating symmetric from nonsymmetric positively curved manifolds?  

**Refined Subproblem:**  
Find explicit examples (including known ones such as homogeneous flag manifolds, Eschenburg spaces, and Bazaikin spaces) and compute \((c_- + c_+)\) for each.  
Do any non‑symmetric positively curved cohomogeneity‑one manifolds satisfy \(c_- + c_+ < n/2\)?  
If none do, is there a geometric reason tied to connectedness or curvature pinching that forbids such configurations?

Why is it a "good" problem:  
This problem evolves the symmetry‑rigidity program toward a **quantitative cohomogeneity–codimension duality**, grounded in the structure of Wilking’s connectedness lemma and empirical properties of known examples.  
It no longer assumes artificial total‑geodesy or roundness—it instead challenges one to understand whether the *combined smallness of orbit codimensions* forces **topological periodicity** or **symmetric rigidity**.  

It thereby integrates three major strands:

- **Riemannian geometry:** curvature and connectedness estimates in positive curvature;
- **Transformation groups:** the detailed orbit‑type stratification of cohomogeneity‑one actions;
- **Topology:** cohomological periodicity and classification via rational homotopy.  

The question is simple to state and falsifiable through computations on known examples, yet answering it would require deep geometric insight into how curvature restrictions propagate through orbit dimensions. By seeking a precise threshold function \(f(n)\), it opens a new quantitative rigidity principle linking curvature, symmetry, and topology—honoring the “narrow entrance, profound depth” ideal of a truly good research‑level problem.